Gangguan pendengaran (\textit{hearing loss}) merupakan salah satu masalah kesehatan yang dapat memengaruhi kualitas hidup serta proses tumbuh kembang seseorang, terutama pada bayi dan anak usia dini. Deteksi dini sangat penting untuk mencegah keterlambatan perkembangan komunikasi, bahasa, dan kemampuan sosial. Namun, alat pemeriksaan gangguan pendengaran yang tersedia di pasaran umumnya memiliki harga relatif mahal, membutuhkan tenaga medis terlatih, serta belum mudah diakses oleh masyarakat luas. Permasalahan tersebut mendorong perlunya pengembangan solusi skrining gangguan pendengaran yang lebih ekonomis, praktis, dan mudah digunakan oleh siapa saja.

Pada penelitian ini dirancang sebuah alat \textit{wearable} untuk \textit{screening hearing loss} berbasis Internet of Things (IoT) yang terintegrasi dengan website monitoring dan pengolahan data. Sistem dirancang agar mampu melakukan pemeriksaan awal dalam waktu singkat dengan prosedur penggunaan yang sederhana sehingga dapat digunakan secara mandiri oleh masyarakat. Selain itu, sistem direncanakan mendukung implementasi \textit{machine learning} untuk membantu analisis hasil pemeriksaan dan meningkatkan akurasi identifikasi potensi gangguan pendengaran. Luaran utama penelitian berupa prototipe alat siap pakai beserta platform website yang dapat digunakan untuk menampilkan hasil pemeriksaan dan mendukung proses monitoring data pengguna.

Tahapan penelitian meliputi studi pustaka, identifikasi kebutuhan sistem, perancangan perangkat keras dan perangkat lunak, integrasi sensor dan komunikasi data, serta pengembangan antarmuka website. Sistem juga akan dirancang dengan mempertimbangkan aspek biaya produksi agar lebih ekonomis dibandingkan alat sejenis di pasaran. Hasil penelitian ini diharapkan dapat menjadi solusi awal untuk meningkatkan akses masyarakat terhadap layanan skrining gangguan pendengaran, khususnya bagi bayi dan anak-anak, sehingga proses deteksi dini dapat dilakukan lebih cepat dan mendukung tumbuh kembang yang lebih optimal.
